% !PRINTED 551
\weberpage{551}
In der Differenz $2S_0 - S$ kommen nun genau dieselben
Glieder vor wie in $S$, nur erhalten die, in denen $x+y$ ungerade
ist, das negative Zeichen, und wir können daher auch setzen:
\weq{12}{\begin{aligned}
\frac{-2\pi}{\beta}\log f(\omega_1)\,f(\omega_2) &= \mathop{\mathrm{Lim}}\limits_{s=1} \sum \frac{(-1)^{x+y}}{[(x-\alpha y)^2+\beta^2y^2]^s},\\
\frac{-2\pi}{\beta}\log f_1(\omega_1)\,f_1(\omega_2) &= \mathop{\mathrm{Lim}}\limits_{s=1} \sum \frac{(-1)^{x}}{[(x-\alpha y)^2+\beta^2y^2]^s},\\
\frac{-2\pi}{\beta}\log f_2(\omega_1)\,f_2(\omega_2) &= \mathop{\mathrm{Lim}}\limits_{s=1} \sum \frac{(-1)^{y}}{[(x-\alpha y)^2+\beta^2y^2]^s},
\end{aligned}}
worin nun $x, y$ \gesperrt{alle} ganzzahligen Wertpaare mit Ausnahme von
$0,0$ annehmen. Diese drei Formeln lassen sich nach § 34, (13), (14)
aus einer von ihnen ableiten, z.\,B. indem man in der zweiten
\weq{13}{\begin{array}{rcll}
\omega_1,\ \omega_2 & \text{durch} & \omega_1+1, & \omega_2-1,\\
& \text{und durch} & \dfrac{-1}{\omega_1}, & \dfrac{-1}{\omega_2},\\
\text{also} &&&\\
\alpha,\ \beta & \text{durch} & \alpha+1, & \beta,\\
& \text{und durch} & \dfrac{-\alpha}{\alpha^2+\beta^2}, & \dfrac{\beta}{\alpha^2+\beta^2}
\end{array}}
ersetzt.

\begin{center}\textbf{\S\,146. Ein Satz über Reihenkonvergenz.}\end{center}

Wollte man in den Ausdrücken (12), § 145 die Zeichen Lim
und $\Sigma$ miteinander vertauschen, also ohne weiteres unter dem
Summenzeichen $s=1$ setzen, so würde man keine unbedingt
konvergente Reihen erhalten, und es muß also untersucht werden,
in welchem Sinne diese Formeln dann noch gültig bleiben.

Um diese Untersuchung durchzuführen, wollen wir zunächst
einen allgemeinen Satz aus der Reihenlehre ableiten, der eine
Verschärfung des Satzes Bd.\,II, § 196, 3. ist. Es sei
\weq{1}{\mu_1 \leqq \mu_2 \leqq \mu_3 \leqq \mu_4 \ldots}
eine Reihe von \gesperrt{unendlich vielen} positiven Zahlen, und $Z(t)$
bedeute die Anzahl dieser Zahlen, die nicht größer als $t$ sind.
Dieses $Z(t)$ soll für jedes $t$ einen endlichen Wert haben, woraus
dann folgt, daß die $\mu_n$ mit $n$ ins unendliche wachsen. Wir
wollen aber noch weiter voraussetzen, daß
\weq{2}{\frac{Z(t)}{t} = a + \frac{\theta}{\sqrt{t}}}
%CONT>>
